% =============================================================================
% emqa_preamble.tex -- shared Beamer preamble of Energy Markets Quantitative Analysis (EMQA)
% Used by every lecture: lectures/Lk/lecturek.tex does
%   \documentclass[10pt,aspectratio=169,t]{beamer}
%   % =============================================================================
% emqa_preamble.tex -- shared Beamer preamble of Energy Markets Quantitative Analysis (EMQA)
% Used by every lecture: lectures/Lk/lecturek.tex does
%   \documentclass[10pt,aspectratio=169,t]{beamer}
%   % =============================================================================
% emqa_preamble.tex -- shared Beamer preamble of Energy Markets Quantitative Analysis (EMQA)
% Used by every lecture: lectures/Lk/lecturek.tex does
%   \documentclass[10pt,aspectratio=169,t]{beamer}
%   % =============================================================================
% emqa_preamble.tex -- shared Beamer preamble of Energy Markets Quantitative Analysis (EMQA)
% Used by every lecture: lectures/Lk/lecturek.tex does
%   \documentclass[10pt,aspectratio=169,t]{beamer}
%   \input{../common/emqa_preamble}
%   \emqalecture{k}{Lecture title}
% Logos live in lectures/common/logos/, charts in lectures/Lk/charts/.
% Do NOT add \hfuzz / \vfuzz / \hbadness here: overfull boxes must be fixed, not hidden.
% =============================================================================

\usetheme{default}
\setbeamersize{text margin left=7mm, text margin right=7mm}

% ---------------------------------------------------------------- packages
\usepackage[utf8]{inputenc}
\usepackage[T1]{fontenc}
\usepackage{lmodern}
\usepackage{amsmath, amssymb}
\usepackage{mathtools}
\usepackage{bm}
\usepackage{tikz}
\usetikzlibrary{arrows.meta, positioning, calc, shapes.misc}
\usepackage{booktabs}
\usepackage{array}
\usepackage{graphicx}
\usepackage{colortbl}
\usepackage{xcolor}
\usepackage{multicol}
\hypersetup{colorlinks=true, linkcolor=MainBlue, urlcolor=MainBlue, citecolor=MainBlue}
\graphicspath{{../common/logos/}}   % charts are included with their language folder: charts/en/, charts/ro/

% ---------------------------------------------------------------- colours
\definecolor{MainBlue}{RGB}{26, 58, 110}
\definecolor{IDAred}{RGB}{205, 0, 0}
\definecolor{DarkGray}{RGB}{51, 51, 51}
\definecolor{FooterGray}{RGB}{100, 100, 100}
\definecolor{BlockGray}{RGB}{236, 238, 242}
\definecolor{KeynoteGray}{RGB}{218, 218, 218}
\definecolor{Crimson}{RGB}{190, 40, 50}
\definecolor{Forest}{RGB}{30, 110, 50}
\definecolor{Bitter}{RGB}{214, 150, 20}
% chart palette (same hex values as lectures/common/emqa_plot.py)
\definecolor{SerRO}{HTML}{2A78D6}
\definecolor{SerDE}{HTML}{EB6834}
\definecolor{SerHU}{HTML}{1BAF7A}
\definecolor{SerBG}{HTML}{EDA100}

\setbeamertemplate{background}{%
  \begin{tikzpicture}[remember picture, overlay]
    \shade[shading=axis, shading angle=315, top color=white, bottom color=KeynoteGray]
      (current page.south west) rectangle (current page.north east);
  \end{tikzpicture}}
\setbeamercolor{background canvas}{bg=}
\setbeamercolor{structure}{fg=MainBlue}
\setbeamercolor{normal text}{fg=DarkGray}
\setbeamercolor{title}{fg=IDAred}
\setbeamercolor{frametitle}{fg=IDAred, bg=}
\setbeamercolor{block title}{bg=MainBlue, fg=white}
\setbeamercolor{block body}{bg=BlockGray, fg=DarkGray}
\setbeamercolor{block title alerted}{bg=Crimson, fg=white}
\setbeamercolor{block body alerted}{bg=Crimson!7, fg=DarkGray}
\setbeamercolor{block title example}{bg=Forest, fg=white}
\setbeamercolor{block body example}{bg=Forest!7, fg=DarkGray}
\setbeamercolor{item}{fg=MainBlue}
\setbeamerfont{frametitle}{size=\large, series=\bfseries}
\setbeamerfont{block title}{size=\small, series=\bfseries}
\setbeamertemplate{blocks}[rounded][shadow=false]
\setbeamertemplate{itemize item}{\color{MainBlue}$\boxdot$}
\setbeamertemplate{itemize subitem}{\color{MainBlue}\scriptsize$\blacktriangleright$}
% compact two-level lists: level 1 = main point, level 2 = details (no text outside bullets)
\setbeamerfont{itemize/enumerate subbody}{size=\footnotesize}
\setlength{\leftmargini}{1.1em}
\setlength{\leftmarginii}{1.1em}
\setbeamertemplate{itemize/enumerate body begin}{\setlength{\topsep}{0pt}\setlength{\partopsep}{0pt}\setlength{\parskip}{0pt}}
\setbeamertemplate{itemize/enumerate subbody begin}{\setlength{\topsep}{0pt}\setlength{\partopsep}{0pt}\setlength{\itemsep}{0pt}\vspace{-0.4ex}}

\setbeamertemplate{enumerate item}{\color{MainBlue}\insertenumlabel.}
\setlength{\leftmargini}{1.3em}
\setlength{\leftmarginii}{1.2em}
\setbeamertemplate{navigation symbols}{}
\setbeamertemplate{section in toc}{\color{MainBlue}$\blacktriangleright$~\inserttocsection}
\setbeamertemplate{caption}{\raggedright\insertcaption\par}
% frame subtitle = precise citation line under the claim title (Weron style)
\setbeamerfont{framesubtitle}{size=\footnotesize}
\setbeamercolor{framesubtitle}{fg=FooterGray}
% outline that reappears at every section with the current section highlighted
\setbeamertemplate{section in toc shaded}[default][35]
\setbeamertemplate{subsection in toc shaded}[default][35]
\setbeamertemplate{subsection in toc}{\leavevmode\leftskip=2.2em\color{DarkGray}\small\inserttocsubsection\par}
\AtBeginSection[]{{\begin{frame}{\txtoutline}\vspace{2mm}\small
  \tableofcontents[currentsection,hideothersubsections]\end{frame}}}

% ---------------------------------------------------------------- head and foot
\setbeamertemplate{headline}{%
  \vskip10pt%
  \hbox to \paperwidth{\hskip0.5cm%
    {\small\color{FooterGray}\renewcommand{\hyperlink}[2]{##2}\insertsectionhead}%
    \hfill\textcolor{FooterGray}{\small\insertframenumber}\hskip0.5cm}%
  \vskip4pt{\color{FooterGray}\hrule height 0.4pt}}
% footer as in the other courses (MFM, SFM): course name on the left, small sparkline on the right
\setbeamertemplate{footline}{%
  {\color{FooterGray}\hrule height 0.4pt}\vskip4pt%
  \hbox to \paperwidth{\hskip0.5cm%
    \textcolor{FooterGray}{\small\txtcourse}%
    \hfill\raisebox{-0.1em}{\begin{tikzpicture}[x=0.08em, y=0.08em, line width=0.4pt]
      \draw[FooterGray] (0,3) -- (1,4) -- (2,3.5) -- (3,5) -- (4,4) -- (5,6) -- (6,5.5) -- (7,4) -- (8,5) -- (9,7) -- (10,6) -- (11,5) -- (12,6.5) -- (13,8) -- (14,7) -- (15,6) -- (16,7.5) -- (17,9) -- (18,8) -- (19,7) -- (20,8.5) -- (21,10) -- (22,9) -- (23,8) -- (24,9.5);
    \end{tikzpicture}}\hskip0.5cm}%
  \vskip6pt}

% ---------------------------------------------------------------- title page
\defbeamertemplate*{title page}{emqa}[1][]{%
  \vspace{0.15cm}
  \begin{center}
    \href{https://www.ase.ro}{\includegraphics[height=0.9cm]{ase_logo.png}}\hspace{0.3cm}%
    \href{https://ipe.ro/new}{\includegraphics[height=0.9cm]{acad_logo.png}}\hspace{0.3cm}%
    \href{https://theida.net}{\includegraphics[height=0.9cm]{ida_logo.png}}\hspace{0.3cm}%
    \href{https://blockchain-research-center.com}{\includegraphics[height=0.9cm]{brc_logo.png}}\hspace{0.3cm}%
    \href{https://www.ai4efin.ase.ro}{\includegraphics[height=0.9cm]{ai4efin_logo.png}}\hspace{0.3cm}%
    \href{https://www.digital-finance-msca.com}{\includegraphics[height=0.9cm]{msca_logo.png}}%
  \end{center}
  \vspace{0.45cm}
  \begin{center}
    \begin{minipage}{0.1\textwidth}\centering
      \href{https://quantlet.com}{\includegraphics[height=0.9cm]{ql_logo.png}}\end{minipage}%
    \begin{minipage}{0.78\textwidth}\centering
      {\Large\bfseries\usebeamercolor[fg]{title}\inserttitle}\par\vspace{0.25cm}
      {\large\usebeamercolor[fg]{title}\insertsubtitle}\end{minipage}%
    \begin{minipage}{0.1\textwidth}\centering
      \href{https://quantinar.com}{\includegraphics[height=0.9cm]{qr_logo.png}}\end{minipage}
  \end{center}
  \vfill
  \hspace{0.4cm}{\bfseries\insertauthor}\par\vspace{0.12cm}
  \hspace{0.4cm}\begin{minipage}[t]{0.92\textwidth}\raggedright\footnotesize\insertinstitute\end{minipage}
  \par\vspace{0.4cm}}

% ---------------------------------------------------------------- lecture metadata
\newcommand{\emqashort}{}
\newcommand{\emqanb}{notebooks}
\newcommand{\emqarepo}{https://github.com/danpele/EMQA/tree/main}
% language-dependent labels (English default; \emqaromanian switches to Romanian, call it before \emqalecture)
\newcommand{\txtcourse}{Energy Markets Quantitative Analysis}
\newcommand{\txtlecture}{Lecture}
\newcommand{\txtoutline}{Outline}
\newcommand{\txtintuition}{In one sentence:}
\newcommand{\txtjob}{In practice:}
\newcommand{\txttry}{Experiment yourself:}
\newcommand{\txtwhere}{where}
\newcommand{\txtsim}{SIMULATED DATA (teaching illustration)}
\newcommand{\txtnbsuffix}{}
\newcommand{\txtinstitute}{Bucharest University of Economic Studies, Department of Statistics and Econometrics\\
    Romanian Academy, Institute for Economic Forecasting\\
    IDA Institute Digital Assets \textbullet\ AI4EFin \textbullet\ MSCA Digital Finance}
\newcommand{\emqaromanian}{%
  \renewcommand{\txtcourse}{Analiza cantitativă a piețelor de energie}%
  \renewcommand{\txtlecture}{Cursul}%
  \renewcommand{\txtoutline}{Cuprins}%
  \renewcommand{\txtintuition}{Pe scurt:}%
  \renewcommand{\txtjob}{În practică:}%
  \renewcommand{\txttry}{Exercițiu practic:}%
  \renewcommand{\txtwhere}{unde}%
  \renewcommand{\txtsim}{DATE SIMULATE (doar pentru ilustrare)}%
  \renewcommand{\txtnbsuffix}{\string_ro}%
  \renewcommand{\txtinstitute}{Academia de Studii Economice din București, Departamentul de Statistică și Econometrie\\
    Academia Română, Institutul de Prognoză Economică\\
    IDA Institute Digital Assets \textbullet\ AI4EFin \textbullet\ MSCA Digital Finance}}
\newcommand{\emqalecture}[2]{%
  \renewcommand{\emqashort}{\txtlecture\ #1}%
  \renewcommand{\emqanb}{notebooks/L#1\string_notebook\txtnbsuffix.ipynb}%
  \title[EMQA]{\txtcourse}%
  \subtitle{\txtlecture\ #1: #2}%
  \author{Prof.\ dr.\ Daniel Traian PELE}%
  \institute{\txtinstitute}%
  \date{}}

% ---------------------------------------------------------------- teaching blocks (MBA clarity rule)
% \intuition{one plain sentence}  ->  formula with every symbol defined  ->  real-data example  ->  \manager{one line}
\newcommand{\intuition}[1]{\begin{itemize}\item {\color{MainBlue}\textbf{\txtintuition}}~#1\end{itemize}\vspace{0.6mm}}
\newcommand{\manager}[1]{\colorbox{Forest!10}{\parbox{\dimexpr\linewidth-2\fboxsep\relax}{\footnotesize
  \begin{itemize}\setlength{\itemsep}{0pt}\item \textbf{\color{Forest}\txtjob}~#1\end{itemize}}}}
% Experiment yourself: pointer to the matching notebook section
\newcommand{\tryit}[1]{{\scriptsize\href{\emqarepo/\emqanb}{%
  \raisebox{-0.15em}{\includegraphics[height=0.85em]{ql_logo.png}}\,\textbf{\txttry}~\color{MainBlue}#1}}}
% bottom row: manager line (left) + experiment pointer (right), pinned to the bottom of the frame
\newcommand{\bottomrow}[2]{\par\vfill\noindent\begin{minipage}[b]{0.72\linewidth}#1\end{minipage}\hfill
  \begin{minipage}[b]{0.26\linewidth}\raggedleft #2\end{minipage}\par\vspace{0.5mm}}
\newcommand{\where}[1]{{\footnotesize\textbf{\txtwhere}~#1\par}}

% ---------------------------------------------------------------- helpers
% chart source line under a figure: \chartsource{Source ...; period ...}
\newcommand{\chartsource}[1]{\par\vspace{1pt}{\scriptsize\color{FooterGray}#1\par}}
% link to code: \quantlet{label}{path inside the repo}
\newcommand{\quantlet}[2]{\par\vspace{1pt}\hfill\href{\emqarepo/#2}{%
  \raisebox{-0.15em}{\includegraphics[height=0.75em]{ql_logo.png}}\textcolor{MainBlue}{\scriptsize\ #1}}}
\newcommand{\simlabel}{\colorbox{Crimson!12}{\scriptsize\bfseries\color{Crimson}\txtsim}}
\newcommand{\E}{\mathbb{E}}
\newcommand{\Var}{\operatorname{Var}}
\newcommand{\Cov}{\operatorname{Cov}}
\newcommand{\Corr}{\operatorname{Corr}}
\newcommand{\asinh}{\operatorname{asinh}}
\newcommand{\iid}{\overset{\text{i.i.d.}}{\sim}}
\newcommand{\EUR}{EUR}

%   \emqalecture{k}{Lecture title}
% Logos live in lectures/common/logos/, charts in lectures/Lk/charts/.
% Do NOT add \hfuzz / \vfuzz / \hbadness here: overfull boxes must be fixed, not hidden.
% =============================================================================

\usetheme{default}
\setbeamersize{text margin left=7mm, text margin right=7mm}

% ---------------------------------------------------------------- packages
\usepackage[utf8]{inputenc}
\usepackage[T1]{fontenc}
\usepackage{lmodern}
\usepackage{amsmath, amssymb}
\usepackage{mathtools}
\usepackage{bm}
\usepackage{tikz}
\usetikzlibrary{arrows.meta, positioning, calc, shapes.misc}
\usepackage{booktabs}
\usepackage{array}
\usepackage{graphicx}
\usepackage{colortbl}
\usepackage{xcolor}
\usepackage{multicol}
\hypersetup{colorlinks=true, linkcolor=MainBlue, urlcolor=MainBlue, citecolor=MainBlue}
\graphicspath{{../common/logos/}}   % charts are included with their language folder: charts/en/, charts/ro/

% ---------------------------------------------------------------- colours
\definecolor{MainBlue}{RGB}{26, 58, 110}
\definecolor{IDAred}{RGB}{205, 0, 0}
\definecolor{DarkGray}{RGB}{51, 51, 51}
\definecolor{FooterGray}{RGB}{100, 100, 100}
\definecolor{BlockGray}{RGB}{236, 238, 242}
\definecolor{KeynoteGray}{RGB}{218, 218, 218}
\definecolor{Crimson}{RGB}{190, 40, 50}
\definecolor{Forest}{RGB}{30, 110, 50}
\definecolor{Bitter}{RGB}{214, 150, 20}
% chart palette (same hex values as lectures/common/emqa_plot.py)
\definecolor{SerRO}{HTML}{2A78D6}
\definecolor{SerDE}{HTML}{EB6834}
\definecolor{SerHU}{HTML}{1BAF7A}
\definecolor{SerBG}{HTML}{EDA100}

\setbeamertemplate{background}{%
  \begin{tikzpicture}[remember picture, overlay]
    \shade[shading=axis, shading angle=315, top color=white, bottom color=KeynoteGray]
      (current page.south west) rectangle (current page.north east);
  \end{tikzpicture}}
\setbeamercolor{background canvas}{bg=}
\setbeamercolor{structure}{fg=MainBlue}
\setbeamercolor{normal text}{fg=DarkGray}
\setbeamercolor{title}{fg=IDAred}
\setbeamercolor{frametitle}{fg=IDAred, bg=}
\setbeamercolor{block title}{bg=MainBlue, fg=white}
\setbeamercolor{block body}{bg=BlockGray, fg=DarkGray}
\setbeamercolor{block title alerted}{bg=Crimson, fg=white}
\setbeamercolor{block body alerted}{bg=Crimson!7, fg=DarkGray}
\setbeamercolor{block title example}{bg=Forest, fg=white}
\setbeamercolor{block body example}{bg=Forest!7, fg=DarkGray}
\setbeamercolor{item}{fg=MainBlue}
\setbeamerfont{frametitle}{size=\large, series=\bfseries}
\setbeamerfont{block title}{size=\small, series=\bfseries}
\setbeamertemplate{blocks}[rounded][shadow=false]
\setbeamertemplate{itemize item}{\color{MainBlue}$\boxdot$}
\setbeamertemplate{itemize subitem}{\color{MainBlue}\scriptsize$\blacktriangleright$}
% compact two-level lists: level 1 = main point, level 2 = details (no text outside bullets)
\setbeamerfont{itemize/enumerate subbody}{size=\footnotesize}
\setlength{\leftmargini}{1.1em}
\setlength{\leftmarginii}{1.1em}
\setbeamertemplate{itemize/enumerate body begin}{\setlength{\topsep}{0pt}\setlength{\partopsep}{0pt}\setlength{\parskip}{0pt}}
\setbeamertemplate{itemize/enumerate subbody begin}{\setlength{\topsep}{0pt}\setlength{\partopsep}{0pt}\setlength{\itemsep}{0pt}\vspace{-0.4ex}}

\setbeamertemplate{enumerate item}{\color{MainBlue}\insertenumlabel.}
\setlength{\leftmargini}{1.3em}
\setlength{\leftmarginii}{1.2em}
\setbeamertemplate{navigation symbols}{}
\setbeamertemplate{section in toc}{\color{MainBlue}$\blacktriangleright$~\inserttocsection}
\setbeamertemplate{caption}{\raggedright\insertcaption\par}
% frame subtitle = precise citation line under the claim title (Weron style)
\setbeamerfont{framesubtitle}{size=\footnotesize}
\setbeamercolor{framesubtitle}{fg=FooterGray}
% outline that reappears at every section with the current section highlighted
\setbeamertemplate{section in toc shaded}[default][35]
\setbeamertemplate{subsection in toc shaded}[default][35]
\setbeamertemplate{subsection in toc}{\leavevmode\leftskip=2.2em\color{DarkGray}\small\inserttocsubsection\par}
\AtBeginSection[]{{\begin{frame}{\txtoutline}\vspace{2mm}\small
  \tableofcontents[currentsection,hideothersubsections]\end{frame}}}

% ---------------------------------------------------------------- head and foot
\setbeamertemplate{headline}{%
  \vskip10pt%
  \hbox to \paperwidth{\hskip0.5cm%
    {\small\color{FooterGray}\renewcommand{\hyperlink}[2]{##2}\insertsectionhead}%
    \hfill\textcolor{FooterGray}{\small\insertframenumber}\hskip0.5cm}%
  \vskip4pt{\color{FooterGray}\hrule height 0.4pt}}
% footer as in the other courses (MFM, SFM): course name on the left, small sparkline on the right
\setbeamertemplate{footline}{%
  {\color{FooterGray}\hrule height 0.4pt}\vskip4pt%
  \hbox to \paperwidth{\hskip0.5cm%
    \textcolor{FooterGray}{\small\txtcourse}%
    \hfill\raisebox{-0.1em}{\begin{tikzpicture}[x=0.08em, y=0.08em, line width=0.4pt]
      \draw[FooterGray] (0,3) -- (1,4) -- (2,3.5) -- (3,5) -- (4,4) -- (5,6) -- (6,5.5) -- (7,4) -- (8,5) -- (9,7) -- (10,6) -- (11,5) -- (12,6.5) -- (13,8) -- (14,7) -- (15,6) -- (16,7.5) -- (17,9) -- (18,8) -- (19,7) -- (20,8.5) -- (21,10) -- (22,9) -- (23,8) -- (24,9.5);
    \end{tikzpicture}}\hskip0.5cm}%
  \vskip6pt}

% ---------------------------------------------------------------- title page
\defbeamertemplate*{title page}{emqa}[1][]{%
  \vspace{0.15cm}
  \begin{center}
    \href{https://www.ase.ro}{\includegraphics[height=0.9cm]{ase_logo.png}}\hspace{0.3cm}%
    \href{https://ipe.ro/new}{\includegraphics[height=0.9cm]{acad_logo.png}}\hspace{0.3cm}%
    \href{https://theida.net}{\includegraphics[height=0.9cm]{ida_logo.png}}\hspace{0.3cm}%
    \href{https://blockchain-research-center.com}{\includegraphics[height=0.9cm]{brc_logo.png}}\hspace{0.3cm}%
    \href{https://www.ai4efin.ase.ro}{\includegraphics[height=0.9cm]{ai4efin_logo.png}}\hspace{0.3cm}%
    \href{https://www.digital-finance-msca.com}{\includegraphics[height=0.9cm]{msca_logo.png}}%
  \end{center}
  \vspace{0.45cm}
  \begin{center}
    \begin{minipage}{0.1\textwidth}\centering
      \href{https://quantlet.com}{\includegraphics[height=0.9cm]{ql_logo.png}}\end{minipage}%
    \begin{minipage}{0.78\textwidth}\centering
      {\Large\bfseries\usebeamercolor[fg]{title}\inserttitle}\par\vspace{0.25cm}
      {\large\usebeamercolor[fg]{title}\insertsubtitle}\end{minipage}%
    \begin{minipage}{0.1\textwidth}\centering
      \href{https://quantinar.com}{\includegraphics[height=0.9cm]{qr_logo.png}}\end{minipage}
  \end{center}
  \vfill
  \hspace{0.4cm}{\bfseries\insertauthor}\par\vspace{0.12cm}
  \hspace{0.4cm}\begin{minipage}[t]{0.92\textwidth}\raggedright\footnotesize\insertinstitute\end{minipage}
  \par\vspace{0.4cm}}

% ---------------------------------------------------------------- lecture metadata
\newcommand{\emqashort}{}
\newcommand{\emqanb}{notebooks}
\newcommand{\emqarepo}{https://github.com/danpele/EMQA/tree/main}
% language-dependent labels (English default; \emqaromanian switches to Romanian, call it before \emqalecture)
\newcommand{\txtcourse}{Energy Markets Quantitative Analysis}
\newcommand{\txtlecture}{Lecture}
\newcommand{\txtoutline}{Outline}
\newcommand{\txtintuition}{In one sentence:}
\newcommand{\txtjob}{In practice:}
\newcommand{\txttry}{Experiment yourself:}
\newcommand{\txtwhere}{where}
\newcommand{\txtsim}{SIMULATED DATA (teaching illustration)}
\newcommand{\txtnbsuffix}{}
\newcommand{\txtinstitute}{Bucharest University of Economic Studies, Department of Statistics and Econometrics\\
    Romanian Academy, Institute for Economic Forecasting\\
    IDA Institute Digital Assets \textbullet\ AI4EFin \textbullet\ MSCA Digital Finance}
\newcommand{\emqaromanian}{%
  \renewcommand{\txtcourse}{Analiza cantitativă a piețelor de energie}%
  \renewcommand{\txtlecture}{Cursul}%
  \renewcommand{\txtoutline}{Cuprins}%
  \renewcommand{\txtintuition}{Pe scurt:}%
  \renewcommand{\txtjob}{În practică:}%
  \renewcommand{\txttry}{Exercițiu practic:}%
  \renewcommand{\txtwhere}{unde}%
  \renewcommand{\txtsim}{DATE SIMULATE (doar pentru ilustrare)}%
  \renewcommand{\txtnbsuffix}{\string_ro}%
  \renewcommand{\txtinstitute}{Academia de Studii Economice din București, Departamentul de Statistică și Econometrie\\
    Academia Română, Institutul de Prognoză Economică\\
    IDA Institute Digital Assets \textbullet\ AI4EFin \textbullet\ MSCA Digital Finance}}
\newcommand{\emqalecture}[2]{%
  \renewcommand{\emqashort}{\txtlecture\ #1}%
  \renewcommand{\emqanb}{notebooks/L#1\string_notebook\txtnbsuffix.ipynb}%
  \title[EMQA]{\txtcourse}%
  \subtitle{\txtlecture\ #1: #2}%
  \author{Prof.\ dr.\ Daniel Traian PELE}%
  \institute{\txtinstitute}%
  \date{}}

% ---------------------------------------------------------------- teaching blocks (MBA clarity rule)
% \intuition{one plain sentence}  ->  formula with every symbol defined  ->  real-data example  ->  \manager{one line}
\newcommand{\intuition}[1]{\begin{itemize}\item {\color{MainBlue}\textbf{\txtintuition}}~#1\end{itemize}\vspace{0.6mm}}
\newcommand{\manager}[1]{\colorbox{Forest!10}{\parbox{\dimexpr\linewidth-2\fboxsep\relax}{\footnotesize
  \begin{itemize}\setlength{\itemsep}{0pt}\item \textbf{\color{Forest}\txtjob}~#1\end{itemize}}}}
% Experiment yourself: pointer to the matching notebook section
\newcommand{\tryit}[1]{{\scriptsize\href{\emqarepo/\emqanb}{%
  \raisebox{-0.15em}{\includegraphics[height=0.85em]{ql_logo.png}}\,\textbf{\txttry}~\color{MainBlue}#1}}}
% bottom row: manager line (left) + experiment pointer (right), pinned to the bottom of the frame
\newcommand{\bottomrow}[2]{\par\vfill\noindent\begin{minipage}[b]{0.72\linewidth}#1\end{minipage}\hfill
  \begin{minipage}[b]{0.26\linewidth}\raggedleft #2\end{minipage}\par\vspace{0.5mm}}
\newcommand{\where}[1]{{\footnotesize\textbf{\txtwhere}~#1\par}}

% ---------------------------------------------------------------- helpers
% chart source line under a figure: \chartsource{Source ...; period ...}
\newcommand{\chartsource}[1]{\par\vspace{1pt}{\scriptsize\color{FooterGray}#1\par}}
% link to code: \quantlet{label}{path inside the repo}
\newcommand{\quantlet}[2]{\par\vspace{1pt}\hfill\href{\emqarepo/#2}{%
  \raisebox{-0.15em}{\includegraphics[height=0.75em]{ql_logo.png}}\textcolor{MainBlue}{\scriptsize\ #1}}}
\newcommand{\simlabel}{\colorbox{Crimson!12}{\scriptsize\bfseries\color{Crimson}\txtsim}}
\newcommand{\E}{\mathbb{E}}
\newcommand{\Var}{\operatorname{Var}}
\newcommand{\Cov}{\operatorname{Cov}}
\newcommand{\Corr}{\operatorname{Corr}}
\newcommand{\asinh}{\operatorname{asinh}}
\newcommand{\iid}{\overset{\text{i.i.d.}}{\sim}}
\newcommand{\EUR}{EUR}

%   \emqalecture{k}{Lecture title}
% Logos live in lectures/common/logos/, charts in lectures/Lk/charts/.
% Do NOT add \hfuzz / \vfuzz / \hbadness here: overfull boxes must be fixed, not hidden.
% =============================================================================

\usetheme{default}
\setbeamersize{text margin left=7mm, text margin right=7mm}

% ---------------------------------------------------------------- packages
\usepackage[utf8]{inputenc}
\usepackage[T1]{fontenc}
\usepackage{lmodern}
\usepackage{amsmath, amssymb}
\usepackage{mathtools}
\usepackage{bm}
\usepackage{tikz}
\usetikzlibrary{arrows.meta, positioning, calc, shapes.misc}
\usepackage{booktabs}
\usepackage{array}
\usepackage{graphicx}
\usepackage{colortbl}
\usepackage{xcolor}
\usepackage{multicol}
\hypersetup{colorlinks=true, linkcolor=MainBlue, urlcolor=MainBlue, citecolor=MainBlue}
\graphicspath{{../common/logos/}}   % charts are included with their language folder: charts/en/, charts/ro/

% ---------------------------------------------------------------- colours
\definecolor{MainBlue}{RGB}{26, 58, 110}
\definecolor{IDAred}{RGB}{205, 0, 0}
\definecolor{DarkGray}{RGB}{51, 51, 51}
\definecolor{FooterGray}{RGB}{100, 100, 100}
\definecolor{BlockGray}{RGB}{236, 238, 242}
\definecolor{KeynoteGray}{RGB}{218, 218, 218}
\definecolor{Crimson}{RGB}{190, 40, 50}
\definecolor{Forest}{RGB}{30, 110, 50}
\definecolor{Bitter}{RGB}{214, 150, 20}
% chart palette (same hex values as lectures/common/emqa_plot.py)
\definecolor{SerRO}{HTML}{2A78D6}
\definecolor{SerDE}{HTML}{EB6834}
\definecolor{SerHU}{HTML}{1BAF7A}
\definecolor{SerBG}{HTML}{EDA100}

\setbeamertemplate{background}{%
  \begin{tikzpicture}[remember picture, overlay]
    \shade[shading=axis, shading angle=315, top color=white, bottom color=KeynoteGray]
      (current page.south west) rectangle (current page.north east);
  \end{tikzpicture}}
\setbeamercolor{background canvas}{bg=}
\setbeamercolor{structure}{fg=MainBlue}
\setbeamercolor{normal text}{fg=DarkGray}
\setbeamercolor{title}{fg=IDAred}
\setbeamercolor{frametitle}{fg=IDAred, bg=}
\setbeamercolor{block title}{bg=MainBlue, fg=white}
\setbeamercolor{block body}{bg=BlockGray, fg=DarkGray}
\setbeamercolor{block title alerted}{bg=Crimson, fg=white}
\setbeamercolor{block body alerted}{bg=Crimson!7, fg=DarkGray}
\setbeamercolor{block title example}{bg=Forest, fg=white}
\setbeamercolor{block body example}{bg=Forest!7, fg=DarkGray}
\setbeamercolor{item}{fg=MainBlue}
\setbeamerfont{frametitle}{size=\large, series=\bfseries}
\setbeamerfont{block title}{size=\small, series=\bfseries}
\setbeamertemplate{blocks}[rounded][shadow=false]
\setbeamertemplate{itemize item}{\color{MainBlue}$\boxdot$}
\setbeamertemplate{itemize subitem}{\color{MainBlue}\scriptsize$\blacktriangleright$}
% compact two-level lists: level 1 = main point, level 2 = details (no text outside bullets)
\setbeamerfont{itemize/enumerate subbody}{size=\footnotesize}
\setlength{\leftmargini}{1.1em}
\setlength{\leftmarginii}{1.1em}
\setbeamertemplate{itemize/enumerate body begin}{\setlength{\topsep}{0pt}\setlength{\partopsep}{0pt}\setlength{\parskip}{0pt}}
\setbeamertemplate{itemize/enumerate subbody begin}{\setlength{\topsep}{0pt}\setlength{\partopsep}{0pt}\setlength{\itemsep}{0pt}\vspace{-0.4ex}}

\setbeamertemplate{enumerate item}{\color{MainBlue}\insertenumlabel.}
\setlength{\leftmargini}{1.3em}
\setlength{\leftmarginii}{1.2em}
\setbeamertemplate{navigation symbols}{}
\setbeamertemplate{section in toc}{\color{MainBlue}$\blacktriangleright$~\inserttocsection}
\setbeamertemplate{caption}{\raggedright\insertcaption\par}
% frame subtitle = precise citation line under the claim title (Weron style)
\setbeamerfont{framesubtitle}{size=\footnotesize}
\setbeamercolor{framesubtitle}{fg=FooterGray}
% outline that reappears at every section with the current section highlighted
\setbeamertemplate{section in toc shaded}[default][35]
\setbeamertemplate{subsection in toc shaded}[default][35]
\setbeamertemplate{subsection in toc}{\leavevmode\leftskip=2.2em\color{DarkGray}\small\inserttocsubsection\par}
\AtBeginSection[]{{\begin{frame}{\txtoutline}\vspace{2mm}\small
  \tableofcontents[currentsection,hideothersubsections]\end{frame}}}

% ---------------------------------------------------------------- head and foot
\setbeamertemplate{headline}{%
  \vskip10pt%
  \hbox to \paperwidth{\hskip0.5cm%
    {\small\color{FooterGray}\renewcommand{\hyperlink}[2]{##2}\insertsectionhead}%
    \hfill\textcolor{FooterGray}{\small\insertframenumber}\hskip0.5cm}%
  \vskip4pt{\color{FooterGray}\hrule height 0.4pt}}
% footer as in the other courses (MFM, SFM): course name on the left, small sparkline on the right
\setbeamertemplate{footline}{%
  {\color{FooterGray}\hrule height 0.4pt}\vskip4pt%
  \hbox to \paperwidth{\hskip0.5cm%
    \textcolor{FooterGray}{\small\txtcourse}%
    \hfill\raisebox{-0.1em}{\begin{tikzpicture}[x=0.08em, y=0.08em, line width=0.4pt]
      \draw[FooterGray] (0,3) -- (1,4) -- (2,3.5) -- (3,5) -- (4,4) -- (5,6) -- (6,5.5) -- (7,4) -- (8,5) -- (9,7) -- (10,6) -- (11,5) -- (12,6.5) -- (13,8) -- (14,7) -- (15,6) -- (16,7.5) -- (17,9) -- (18,8) -- (19,7) -- (20,8.5) -- (21,10) -- (22,9) -- (23,8) -- (24,9.5);
    \end{tikzpicture}}\hskip0.5cm}%
  \vskip6pt}

% ---------------------------------------------------------------- title page
\defbeamertemplate*{title page}{emqa}[1][]{%
  \vspace{0.15cm}
  \begin{center}
    \href{https://www.ase.ro}{\includegraphics[height=0.9cm]{ase_logo.png}}\hspace{0.3cm}%
    \href{https://ipe.ro/new}{\includegraphics[height=0.9cm]{acad_logo.png}}\hspace{0.3cm}%
    \href{https://theida.net}{\includegraphics[height=0.9cm]{ida_logo.png}}\hspace{0.3cm}%
    \href{https://blockchain-research-center.com}{\includegraphics[height=0.9cm]{brc_logo.png}}\hspace{0.3cm}%
    \href{https://www.ai4efin.ase.ro}{\includegraphics[height=0.9cm]{ai4efin_logo.png}}\hspace{0.3cm}%
    \href{https://www.digital-finance-msca.com}{\includegraphics[height=0.9cm]{msca_logo.png}}%
  \end{center}
  \vspace{0.45cm}
  \begin{center}
    \begin{minipage}{0.1\textwidth}\centering
      \href{https://quantlet.com}{\includegraphics[height=0.9cm]{ql_logo.png}}\end{minipage}%
    \begin{minipage}{0.78\textwidth}\centering
      {\Large\bfseries\usebeamercolor[fg]{title}\inserttitle}\par\vspace{0.25cm}
      {\large\usebeamercolor[fg]{title}\insertsubtitle}\end{minipage}%
    \begin{minipage}{0.1\textwidth}\centering
      \href{https://quantinar.com}{\includegraphics[height=0.9cm]{qr_logo.png}}\end{minipage}
  \end{center}
  \vfill
  \hspace{0.4cm}{\bfseries\insertauthor}\par\vspace{0.12cm}
  \hspace{0.4cm}\begin{minipage}[t]{0.92\textwidth}\raggedright\footnotesize\insertinstitute\end{minipage}
  \par\vspace{0.4cm}}

% ---------------------------------------------------------------- lecture metadata
\newcommand{\emqashort}{}
\newcommand{\emqanb}{notebooks}
\newcommand{\emqarepo}{https://github.com/danpele/EMQA/tree/main}
% language-dependent labels (English default; \emqaromanian switches to Romanian, call it before \emqalecture)
\newcommand{\txtcourse}{Energy Markets Quantitative Analysis}
\newcommand{\txtlecture}{Lecture}
\newcommand{\txtoutline}{Outline}
\newcommand{\txtintuition}{In one sentence:}
\newcommand{\txtjob}{In practice:}
\newcommand{\txttry}{Experiment yourself:}
\newcommand{\txtwhere}{where}
\newcommand{\txtsim}{SIMULATED DATA (teaching illustration)}
\newcommand{\txtnbsuffix}{}
\newcommand{\txtinstitute}{Bucharest University of Economic Studies, Department of Statistics and Econometrics\\
    Romanian Academy, Institute for Economic Forecasting\\
    IDA Institute Digital Assets \textbullet\ AI4EFin \textbullet\ MSCA Digital Finance}
\newcommand{\emqaromanian}{%
  \renewcommand{\txtcourse}{Analiza cantitativă a piețelor de energie}%
  \renewcommand{\txtlecture}{Cursul}%
  \renewcommand{\txtoutline}{Cuprins}%
  \renewcommand{\txtintuition}{Pe scurt:}%
  \renewcommand{\txtjob}{În practică:}%
  \renewcommand{\txttry}{Exercițiu practic:}%
  \renewcommand{\txtwhere}{unde}%
  \renewcommand{\txtsim}{DATE SIMULATE (doar pentru ilustrare)}%
  \renewcommand{\txtnbsuffix}{\string_ro}%
  \renewcommand{\txtinstitute}{Academia de Studii Economice din București, Departamentul de Statistică și Econometrie\\
    Academia Română, Institutul de Prognoză Economică\\
    IDA Institute Digital Assets \textbullet\ AI4EFin \textbullet\ MSCA Digital Finance}}
\newcommand{\emqalecture}[2]{%
  \renewcommand{\emqashort}{\txtlecture\ #1}%
  \renewcommand{\emqanb}{notebooks/L#1\string_notebook\txtnbsuffix.ipynb}%
  \title[EMQA]{\txtcourse}%
  \subtitle{\txtlecture\ #1: #2}%
  \author{Prof.\ dr.\ Daniel Traian PELE}%
  \institute{\txtinstitute}%
  \date{}}

% ---------------------------------------------------------------- teaching blocks (MBA clarity rule)
% \intuition{one plain sentence}  ->  formula with every symbol defined  ->  real-data example  ->  \manager{one line}
\newcommand{\intuition}[1]{\begin{itemize}\item {\color{MainBlue}\textbf{\txtintuition}}~#1\end{itemize}\vspace{0.6mm}}
\newcommand{\manager}[1]{\colorbox{Forest!10}{\parbox{\dimexpr\linewidth-2\fboxsep\relax}{\footnotesize
  \begin{itemize}\setlength{\itemsep}{0pt}\item \textbf{\color{Forest}\txtjob}~#1\end{itemize}}}}
% Experiment yourself: pointer to the matching notebook section
\newcommand{\tryit}[1]{{\scriptsize\href{\emqarepo/\emqanb}{%
  \raisebox{-0.15em}{\includegraphics[height=0.85em]{ql_logo.png}}\,\textbf{\txttry}~\color{MainBlue}#1}}}
% bottom row: manager line (left) + experiment pointer (right), pinned to the bottom of the frame
\newcommand{\bottomrow}[2]{\par\vfill\noindent\begin{minipage}[b]{0.72\linewidth}#1\end{minipage}\hfill
  \begin{minipage}[b]{0.26\linewidth}\raggedleft #2\end{minipage}\par\vspace{0.5mm}}
\newcommand{\where}[1]{{\footnotesize\textbf{\txtwhere}~#1\par}}

% ---------------------------------------------------------------- helpers
% chart source line under a figure: \chartsource{Source ...; period ...}
\newcommand{\chartsource}[1]{\par\vspace{1pt}{\scriptsize\color{FooterGray}#1\par}}
% link to code: \quantlet{label}{path inside the repo}
\newcommand{\quantlet}[2]{\par\vspace{1pt}\hfill\href{\emqarepo/#2}{%
  \raisebox{-0.15em}{\includegraphics[height=0.75em]{ql_logo.png}}\textcolor{MainBlue}{\scriptsize\ #1}}}
\newcommand{\simlabel}{\colorbox{Crimson!12}{\scriptsize\bfseries\color{Crimson}\txtsim}}
\newcommand{\E}{\mathbb{E}}
\newcommand{\Var}{\operatorname{Var}}
\newcommand{\Cov}{\operatorname{Cov}}
\newcommand{\Corr}{\operatorname{Corr}}
\newcommand{\asinh}{\operatorname{asinh}}
\newcommand{\iid}{\overset{\text{i.i.d.}}{\sim}}
\newcommand{\EUR}{EUR}

%   \emqalecture{k}{Lecture title}
% Logos live in lectures/common/logos/, charts in lectures/Lk/charts/.
% Do NOT add \hfuzz / \vfuzz / \hbadness here: overfull boxes must be fixed, not hidden.
% =============================================================================

\usetheme{default}
\setbeamersize{text margin left=7mm, text margin right=7mm}

% ---------------------------------------------------------------- packages
\usepackage[utf8]{inputenc}
\usepackage[T1]{fontenc}
\usepackage{lmodern}
\usepackage{amsmath, amssymb}
\usepackage{mathtools}
\usepackage{bm}
\usepackage{tikz}
\usetikzlibrary{arrows.meta, positioning, calc, shapes.misc}
\usepackage{booktabs}
\usepackage{array}
\usepackage{graphicx}
\usepackage{colortbl}
\usepackage{xcolor}
\usepackage{multicol}
\hypersetup{colorlinks=true, linkcolor=MainBlue, urlcolor=MainBlue, citecolor=MainBlue}
\graphicspath{{../common/logos/}}   % charts are included with their language folder: charts/en/, charts/ro/

% ---------------------------------------------------------------- colours
\definecolor{MainBlue}{RGB}{26, 58, 110}
\definecolor{IDAred}{RGB}{205, 0, 0}
\definecolor{DarkGray}{RGB}{51, 51, 51}
\definecolor{FooterGray}{RGB}{100, 100, 100}
\definecolor{BlockGray}{RGB}{236, 238, 242}
\definecolor{KeynoteGray}{RGB}{218, 218, 218}
\definecolor{Crimson}{RGB}{190, 40, 50}
\definecolor{Forest}{RGB}{30, 110, 50}
\definecolor{Bitter}{RGB}{214, 150, 20}
% chart palette (same hex values as lectures/common/emqa_plot.py)
\definecolor{SerRO}{HTML}{2A78D6}
\definecolor{SerDE}{HTML}{EB6834}
\definecolor{SerHU}{HTML}{1BAF7A}
\definecolor{SerBG}{HTML}{EDA100}

\setbeamertemplate{background}{%
  \begin{tikzpicture}[remember picture, overlay]
    \shade[shading=axis, shading angle=315, top color=white, bottom color=KeynoteGray]
      (current page.south west) rectangle (current page.north east);
  \end{tikzpicture}}
\setbeamercolor{background canvas}{bg=}
\setbeamercolor{structure}{fg=MainBlue}
\setbeamercolor{normal text}{fg=DarkGray}
\setbeamercolor{title}{fg=IDAred}
\setbeamercolor{frametitle}{fg=IDAred, bg=}
\setbeamercolor{block title}{bg=MainBlue, fg=white}
\setbeamercolor{block body}{bg=BlockGray, fg=DarkGray}
\setbeamercolor{block title alerted}{bg=Crimson, fg=white}
\setbeamercolor{block body alerted}{bg=Crimson!7, fg=DarkGray}
\setbeamercolor{block title example}{bg=Forest, fg=white}
\setbeamercolor{block body example}{bg=Forest!7, fg=DarkGray}
\setbeamercolor{item}{fg=MainBlue}
\setbeamerfont{frametitle}{size=\large, series=\bfseries}
\setbeamerfont{block title}{size=\small, series=\bfseries}
\setbeamertemplate{blocks}[rounded][shadow=false]
\setbeamertemplate{itemize item}{\color{MainBlue}$\boxdot$}
\setbeamertemplate{itemize subitem}{\color{MainBlue}\scriptsize$\blacktriangleright$}
% compact two-level lists: level 1 = main point, level 2 = details (no text outside bullets)
\setbeamerfont{itemize/enumerate subbody}{size=\footnotesize}
\setlength{\leftmargini}{1.1em}
\setlength{\leftmarginii}{1.1em}
\setbeamertemplate{itemize/enumerate body begin}{\setlength{\topsep}{0pt}\setlength{\partopsep}{0pt}\setlength{\parskip}{0pt}}
\setbeamertemplate{itemize/enumerate subbody begin}{\setlength{\topsep}{0pt}\setlength{\partopsep}{0pt}\setlength{\itemsep}{0pt}\vspace{-0.4ex}}

\setbeamertemplate{enumerate item}{\color{MainBlue}\insertenumlabel.}
\setlength{\leftmargini}{1.3em}
\setlength{\leftmarginii}{1.2em}
\setbeamertemplate{navigation symbols}{}
\setbeamertemplate{section in toc}{\color{MainBlue}$\blacktriangleright$~\inserttocsection}
\setbeamertemplate{caption}{\raggedright\insertcaption\par}
% frame subtitle = precise citation line under the claim title (Weron style)
\setbeamerfont{framesubtitle}{size=\footnotesize}
\setbeamercolor{framesubtitle}{fg=FooterGray}
% outline that reappears at every section with the current section highlighted
\setbeamertemplate{section in toc shaded}[default][35]
\setbeamertemplate{subsection in toc shaded}[default][35]
\setbeamertemplate{subsection in toc}{\leavevmode\leftskip=2.2em\color{DarkGray}\small\inserttocsubsection\par}
\AtBeginSection[]{{\begin{frame}{\txtoutline}\vspace{2mm}\small
  \tableofcontents[currentsection,hideothersubsections]\end{frame}}}

% ---------------------------------------------------------------- head and foot
\setbeamertemplate{headline}{%
  \vskip10pt%
  \hbox to \paperwidth{\hskip0.5cm%
    {\small\color{FooterGray}\renewcommand{\hyperlink}[2]{##2}\insertsectionhead}%
    \hfill\textcolor{FooterGray}{\small\insertframenumber}\hskip0.5cm}%
  \vskip4pt{\color{FooterGray}\hrule height 0.4pt}}
% footer as in the other courses (MFM, SFM): course name on the left, small sparkline on the right
\setbeamertemplate{footline}{%
  {\color{FooterGray}\hrule height 0.4pt}\vskip4pt%
  \hbox to \paperwidth{\hskip0.5cm%
    \textcolor{FooterGray}{\small\txtcourse}%
    \hfill\raisebox{-0.1em}{\begin{tikzpicture}[x=0.08em, y=0.08em, line width=0.4pt]
      \draw[FooterGray] (0,3) -- (1,4) -- (2,3.5) -- (3,5) -- (4,4) -- (5,6) -- (6,5.5) -- (7,4) -- (8,5) -- (9,7) -- (10,6) -- (11,5) -- (12,6.5) -- (13,8) -- (14,7) -- (15,6) -- (16,7.5) -- (17,9) -- (18,8) -- (19,7) -- (20,8.5) -- (21,10) -- (22,9) -- (23,8) -- (24,9.5);
    \end{tikzpicture}}\hskip0.5cm}%
  \vskip6pt}

% ---------------------------------------------------------------- title page
\defbeamertemplate*{title page}{emqa}[1][]{%
  \vspace{0.15cm}
  \begin{center}
    \href{https://www.ase.ro}{\includegraphics[height=0.9cm]{ase_logo.png}}\hspace{0.3cm}%
    \href{https://ipe.ro/new}{\includegraphics[height=0.9cm]{acad_logo.png}}\hspace{0.3cm}%
    \href{https://theida.net}{\includegraphics[height=0.9cm]{ida_logo.png}}\hspace{0.3cm}%
    \href{https://blockchain-research-center.com}{\includegraphics[height=0.9cm]{brc_logo.png}}\hspace{0.3cm}%
    \href{https://www.ai4efin.ase.ro}{\includegraphics[height=0.9cm]{ai4efin_logo.png}}\hspace{0.3cm}%
    \href{https://www.digital-finance-msca.com}{\includegraphics[height=0.9cm]{msca_logo.png}}%
  \end{center}
  \vspace{0.45cm}
  \begin{center}
    \begin{minipage}{0.1\textwidth}\centering
      \href{https://quantlet.com}{\includegraphics[height=0.9cm]{ql_logo.png}}\end{minipage}%
    \begin{minipage}{0.78\textwidth}\centering
      {\Large\bfseries\usebeamercolor[fg]{title}\inserttitle}\par\vspace{0.25cm}
      {\large\usebeamercolor[fg]{title}\insertsubtitle}\end{minipage}%
    \begin{minipage}{0.1\textwidth}\centering
      \href{https://quantinar.com}{\includegraphics[height=0.9cm]{qr_logo.png}}\end{minipage}
  \end{center}
  \vfill
  \hspace{0.4cm}{\bfseries\insertauthor}\par\vspace{0.12cm}
  \hspace{0.4cm}\begin{minipage}[t]{0.92\textwidth}\raggedright\footnotesize\insertinstitute\end{minipage}
  \par\vspace{0.4cm}}

% ---------------------------------------------------------------- lecture metadata
\newcommand{\emqashort}{}
\newcommand{\emqanb}{notebooks}
\newcommand{\emqarepo}{https://github.com/danpele/EMQA/tree/main}
% language-dependent labels (English default; \emqaromanian switches to Romanian, call it before \emqalecture)
\newcommand{\txtcourse}{Energy Markets Quantitative Analysis}
\newcommand{\txtlecture}{Lecture}
\newcommand{\txtoutline}{Outline}
\newcommand{\txtintuition}{In one sentence:}
\newcommand{\txtjob}{In practice:}
\newcommand{\txttry}{Experiment yourself:}
\newcommand{\txtwhere}{where}
\newcommand{\txtsim}{SIMULATED DATA (teaching illustration)}
\newcommand{\txtnbsuffix}{}
\newcommand{\txtinstitute}{Bucharest University of Economic Studies, Department of Statistics and Econometrics\\
    Romanian Academy, Institute for Economic Forecasting\\
    IDA Institute Digital Assets \textbullet\ AI4EFin \textbullet\ MSCA Digital Finance}
\newcommand{\emqaromanian}{%
  \renewcommand{\txtcourse}{Analiza cantitativă a piețelor de energie}%
  \renewcommand{\txtlecture}{Cursul}%
  \renewcommand{\txtoutline}{Cuprins}%
  \renewcommand{\txtintuition}{Pe scurt:}%
  \renewcommand{\txtjob}{În practică:}%
  \renewcommand{\txttry}{Exercițiu practic:}%
  \renewcommand{\txtwhere}{unde}%
  \renewcommand{\txtsim}{DATE SIMULATE (doar pentru ilustrare)}%
  \renewcommand{\txtnbsuffix}{\string_ro}%
  \renewcommand{\txtinstitute}{Academia de Studii Economice din București, Departamentul de Statistică și Econometrie\\
    Academia Română, Institutul de Prognoză Economică\\
    IDA Institute Digital Assets \textbullet\ AI4EFin \textbullet\ MSCA Digital Finance}}
\newcommand{\emqalecture}[2]{%
  \renewcommand{\emqashort}{\txtlecture\ #1}%
  \renewcommand{\emqanb}{notebooks/L#1\string_notebook\txtnbsuffix.ipynb}%
  \title[EMQA]{\txtcourse}%
  \subtitle{\txtlecture\ #1: #2}%
  \author{Prof.\ dr.\ Daniel Traian PELE}%
  \institute{\txtinstitute}%
  \date{}}

% ---------------------------------------------------------------- teaching blocks (MBA clarity rule)
% \intuition{one plain sentence}  ->  formula with every symbol defined  ->  real-data example  ->  \manager{one line}
\newcommand{\intuition}[1]{\begin{itemize}\item {\color{MainBlue}\textbf{\txtintuition}}~#1\end{itemize}\vspace{0.6mm}}
\newcommand{\manager}[1]{\colorbox{Forest!10}{\parbox{\dimexpr\linewidth-2\fboxsep\relax}{\footnotesize
  \begin{itemize}\setlength{\itemsep}{0pt}\item \textbf{\color{Forest}\txtjob}~#1\end{itemize}}}}
% Experiment yourself: pointer to the matching notebook section
\newcommand{\tryit}[1]{{\scriptsize\href{\emqarepo/\emqanb}{%
  \raisebox{-0.15em}{\includegraphics[height=0.85em]{ql_logo.png}}\,\textbf{\txttry}~\color{MainBlue}#1}}}
% bottom row: manager line (left) + experiment pointer (right), pinned to the bottom of the frame
\newcommand{\bottomrow}[2]{\par\vfill\noindent\begin{minipage}[b]{0.72\linewidth}#1\end{minipage}\hfill
  \begin{minipage}[b]{0.26\linewidth}\raggedleft #2\end{minipage}\par\vspace{0.5mm}}
\newcommand{\where}[1]{{\footnotesize\textbf{\txtwhere}~#1\par}}

% ---------------------------------------------------------------- helpers
% chart source line under a figure: \chartsource{Source ...; period ...}
\newcommand{\chartsource}[1]{\par\vspace{1pt}{\scriptsize\color{FooterGray}#1\par}}
% link to code: \quantlet{label}{path inside the repo}
\newcommand{\quantlet}[2]{\par\vspace{1pt}\hfill\href{\emqarepo/#2}{%
  \raisebox{-0.15em}{\includegraphics[height=0.75em]{ql_logo.png}}\textcolor{MainBlue}{\scriptsize\ #1}}}
\newcommand{\simlabel}{\colorbox{Crimson!12}{\scriptsize\bfseries\color{Crimson}\txtsim}}
\newcommand{\E}{\mathbb{E}}
\newcommand{\Var}{\operatorname{Var}}
\newcommand{\Cov}{\operatorname{Cov}}
\newcommand{\Corr}{\operatorname{Corr}}
\newcommand{\asinh}{\operatorname{asinh}}
\newcommand{\iid}{\overset{\text{i.i.d.}}{\sim}}
\newcommand{\EUR}{EUR}
